\scanpage{61}

Le résultat qui suit, qui résoud le``problème du cycle invariant'' dans le cas du $H^1$, est essentiellement une reformulation d'un résultat de M.RAYNAUD [\quad]:
\margin{$\times$}

\uline{Théorème} / Soient $S$ un trait strictement local, $f:X\unsure{\longrightarrow|\mathbin{-}} S$ un schéma pro\deleted{\unsure{pre}}\inserted{jectif} sur $S$, régulier, tel que l'on ait $\underline{O}_S\xrightarrow{\sim}f_*(\underline{O}_X)$, (ce qui implique que $X$ est plat sur $S$, et que $X_{\unsure{0|o}}$ est un diviseur sur $X$), \deleted{\unsure{Soient} $\unsure{(D_i)_{i\in J}}$ \unsure{les}}\inserted{$(D_{\unsure{i|m}})_{i\in J}$ la famille des}\begin{tikzpicture}[baseline=-.5ex,x=1mm,y=1mm]\draw (4,3)..controls(5,0)and(0,2)..(0,0)..controls(0,-1)and(.5,-2)..(.6,-3);\end{tikzpicture} composantes irréductibles réduites de $X_{\unsure{0|o}}$, considérées comme diviseurs sur $X$, et posons \deleted{$X$}

\[
(1)\qquad X_{\unsure{0|o}}=\sum_{i\in J}d_iD_i,\qquad\text{où }d_i\geqslant1,
\]

\[
(2)\qquad d=\operatorname{pgcd}(d_i)_{i\in J}.
\]

Soient
\[
(3)\qquad\Delta\simeq\underline{Z}^{J}
\]

la catégorie des diviseurs $D$ sur $X$ tels que $\operatorname{supp}D\subset X_{\unsure{0|o}}$, \deleted{et}

\[
(4)\qquad\deleted{\unsure{Ixx\underline Z}}\ \Gamma\simeq/\underline Z.X_{\unsure{0|o}}\hookrightarrow\operatorname{Pic}(X)
\]

le \deleted{est} sous-groupe de $\operatorname{Pic}(X)$ image de $\Delta$ par l'homomorphisme $D\mapsto\operatorname{cl}(\underline{O}_X(D))$ de $\Delta$ dans $\operatorname{Pic}(X)$, et considérons l'homomorphisme \deleted{compos\unsure{é}}

\[
(5)\qquad w:\Gamma\longrightarrow\operatorname{NS}(X_{\unsure{0|o}})\overset{\mathrm{dfn}}{=}\operatorname{Im}(\operatorname{Pic}(X_{\unsure{0|o}})\longrightarrow\underline{\operatorname{NS}}_{X_{\unsure{0|o}}/k}(k)),
\]

composé des homomorphismes évidents
\diagram{\begin{tikzcd}[ampersand replacement=\&]
\Gamma\arrow[r,hook,"i"{pos=.25}]\&\operatorname{Pic}(X)\arrow[r]\&\operatorname{Pic}(X_{\unsure{0|o}})\arrow[r]\&\operatorname{NS}(X_{\unsure{0|o}}).
\end{tikzcd}}

Soit enfin $d^t$ le pgcd des degrés des 0-cycles sur $X_{\deleted{\unsure{y}}\inserted{\eta}}$. Avec ces notations, on a ce qui suit:

a) L'entier $d^t$ est aussi égal au pgcd des $\mu_i d_i$, où pour tout $i\in J$, $\mu_i$ désigne la multiplicité radicielle de $D_i$ sur $k$ (EGA IV 4.\quad); en particulier \inserted{on a} \deleted{$d$ divise $d^t$, donc est égal à 1 si $d^t$ l'est} \deleted{\inserted{et lui est égal si $k$ est parfait.}} \deleted{\inserted{, \unsure{la réciproque étant} vraie si $k$ est parfait}} \inserted{$d^t=p^h d$, où $p^h=\operatorname{pgcd}(\mu_i)$,}

\margin{\unsure{donc} $h=0$ ie. $d=d_t$ si $k$ est parfait}

\margin{$\times$}

b) L'homomorphisme $w$ est un homomorphisme de $\underline Z$-modules \uline{de type fini dont le noyau est contenu dans le sous-groupe de torsion de} $\Gamma$. Ce dernier est isomorphe canoniquement à $\underline Z/d\underline Z$, avec comme générateur l'image dans $\Gamma$ du cycle $(1/d)X_{\unsure{0|o}}=\sum_{i\in J}(d_i/d)D_i$. Donc $\operatorname{Ker}w$ est cyclique

\margin{\figure{\begin{tikzpicture}[x=.005cm,y=-.005cm]
\draw (185,31)--(186,317)--(187,472)--(187,477)--(187,597)--(187,602)--(189,727)--(191,855)--(195,962)--(201,1087)--(208,1212)--(214,1429);
\draw (198,1392)..controls(196,1432)and(197,1454)..(199,1465)--(236,1464);
\draw (173,1531)..controls(182,1505)and(214,1481)..(236,1464);
\draw (190,1552)--(193,1693);
\draw (0,1430)--(13,1539);
\draw (66,74)--(108,0);
\end{tikzpicture}}}

\margin{60}
